\section{Introdução}

\hspace{1cm} Este trabalho tem como objetivo analisar as diferenças entre a Autenticação Simples (Single-Factor Authentication -- SFA) e a Autenticação Multifator (Multi-Factor Authentication -- MFA), avaliando as suas características, vantagens, limitações e adequação a diferentes contextos de utilização.

\hspace{1cm} Pretende-se compreender em que situações a implementação de MFA é necessária ou recomendada, considerando o nível de risco, a sensibilidade da informação protegida e o perfil dos utilizadores. O trabalho procura também analisar o equilíbrio entre segurança, facilidade de utilização e acessibilidade.

\hspace{1cm} Por fim, serão estudadas algumas das principais tecnologias de autenticação, analisando as duas características e o modo como podem ser utilizadas na implementação de mecanismos de MFA.

\hspace{1cm} O relatório encontra-se organizado em várias secções. Inicialmente, são apresentados os principais tipos de autenticação e os conceitos de SFA e MFA. A seguir, são comparadas ambas as abordagens e analisada a necessidade de recorrer a MFA em função do risco e do perfil dos utilizadores. Por fim, são exploradas diferentes tecnologias de autenticação e apresentados cenários que permitem avaliar a adequação das soluções estudadas a diferentes contextos. 